%!TEX output_directory = ./tmp/
%!TEX program = pdflatex

\documentclass[10pt, a4paper]{article}
\usepackage{tikz}
\usepackage[utf8]{inputenc}
\usepackage[spanish]{babel}
\usepackage{amsmath}
\usepackage{amsfonts}
\usepackage{amssymb}
\usepackage{graphicx}
\usepackage{float}
\usepackage[
    top=2cm, bottom=2cm, left=2cm, right=2cm
]{geometry}

\begin{document}

\title{Student Behaviour Graph Design Document}
\author{
    {José Narciso Robles Durán} \\
    {Universidad Complutense de Madrid} \\
    {Departamento de Informática e Inteligencia Artificial | ActLab}
}
\date{\today}
\maketitle

\tableofcontents
\section{Motivación}
    La herramienta \textit{Didascalias-VR} permite interactuar con estudiantes virtuales dirigidos por comportamientos por inteligencia artificial.\par

    Estos comportamientos son programables en el motor \textit{Unity}. Actualmente, se usa una máquina de estados para escoger y transicionar entre ellos. Los estados, junto con sus transiciones, se reprsentan con \textit{ScriptableObjects}.\par

    Cada estado contiene una lista de transiciones que se evalúan cada frame si eres el estado activo. Si una transición se cumple, se cambia al estado objetivo de esa transición.\par

    Esta implementación tiene como ventaja el razonamiento local para un diseñador de las transiciones para un estado dado. En el inspector de \textit{Unity}, se ve una lista detallada de las referencias a los otros \textit{ScriptableObjects} que representan las transiciones entre dos estados.\par

    Sin embargo, esta implementación tiene como desventaja que no se puede ver el panorama general de los estados y sus transiciones, así como el flujo de la red. Esto hace difícil detectar errores de diseño, como estados sin transiciones incidentes o transiciones que no se cumplen nunca.\par

\section{Propuesta}


\end{document}